\subsection{Alias, Threshold, and Decision Redundancy}
\label{subsec:static-admissibility-alias-threshold-redundancy}

Feature redundancy is not limited to correlation. In the static admissibility
stage, two features may be redundant because they are semantic aliases,
deterministic transformations, near-threshold variants, logically subsuming
indicators, or decision-equivalent measurements. The audit therefore combines
schema-level semantic grouping, functional-equivalence tests,
threshold-subsumption checks, observability-aware threshold-gap analysis, and
conditional decision-gain estimates.

A near-threshold pair such as \texttt{min\_altitude\_ge\_300m} and
\texttt{min\_altitude\_ge\_301m} is not automatically collapsed if the boundary
slice contains class-mixed samples, but if the threshold gap is below declared
measurement uncertainty, retention is labeled candidate-level until ablation or
observability analysis confirms that the distinction is meaningful.

The practical actions are intentionally narrow:
\begin{itemize}
  \item exact duplicates are dropped,
  \item unit or affine aliases are canonicalized,
  \item threshold aliases are collapsed unless the boundary slice changes
        decision evidence,
  \item pair-specific retained thresholds remain candidate-level when the
        observability gap is weak,
  \item decision-redundant features are deprioritized unless they remain useful
        for a named hard class pair.
\end{itemize}
